\mychapter{Introdução}{cap:introducao}
\lhead{INTRODUÇÃO}

Construir a interface de usuário está entre as tarefas mais custosas do desenvolvimento de aplicações móveis, a ponto de consumir mais da metade do esforço total \cite{gao2026guimigrator}. Manter uma versão nativa por plataforma multiplica esse custo, pois exige equipes e bases de código separadas. Por isso, a demanda por soluções eficientes e de baixo custo impulsionou a adoção de frameworks multiplataforma \cite{jost2025state}. Esses frameworks reduzem a duplicação, mas cobram em desempenho e uso de recursos e criam forte dependência da ferramenta \cite{biorn2020empirical, mahmoud2024trans, schultes2025native}. Em uma comparação empírica, Biørn-Hansen et al.\ mediram CPU, memória e tempo de resposta de cinco frameworks e de uma implementação nativa e concluíram que o desempenho multiplataforma é, em geral, inferior ao nativo, sem que nenhum framework seja o pior ou o melhor em todas as métricas \cite{biorn2020empirical}.

Abordagens recentes substituem a camada intermediária por transpilação determinística. O SequalsK compartilha a camada de modelo entre Kotlin e Swift e permite escrever até 86\% do código uma única vez, mas não cobre a interface, que os próprios autores deixam como trabalho futuro \cite{schultes2025native}. O GUIMigrator converte interfaces Android em XML para Jetpack Compose e SwiftUI, com similaridade visual média de 81,9\% e 78,2\%, respectivamente \cite{gao2026guimigrator}. Mahmoud et al.\ convertem aplicações React Native e Xamarin em aplicações nativas, que superaram as originais em memória, armazenamento e velocidade \cite{mahmoud2024trans}. Na literatura triada, os três trabalhos têm Android e iOS como alvo e nenhum gera interfaces nativas de desktop (Windows e macOS) a partir de uma fonte única. No desktop, os estudos encontrados tratam dos desafios de aplicações baseadas em web view \cite{scoccia2021challenges} e da comparação entre Electron, Tauri e PyQt quanto à integração com o sistema, à memória e ao tamanho do instalador \cite{milichiewicz2025comparative}.

Diante dessa lacuna, este trabalho tem como objetivo avaliar se a geração de projetos nativos para cada sistema operacional (SwiftUI no macOS e WinUI no Windows) a partir de uma única fonte produz aplicações com desempenho comparável e alta similaridade visual em relação às escritas diretamente nas ferramentas oficiais de cada plataforma, sem a camada de renderização própria dos frameworks atuais.
